\question[4]
Gib für jede Anweisung an, was ausgegeben wird, und welchen Datentyp der ausgegebene Wert hat.

\begin{lstlisting}
print(7 / 2)
print(7 // 2)
print(int(7.9))
print(int(-2.5))
print(float(4))
print(int('12') + 3)
print('12' + str(3))
print(True + True)
\end{lstlisting}

\begin{solutionbox}{5.5cm}
\begin{lstlisting}[language={}]
3.5    float   (/ liefert immer einen float)
3      int     (ganzzahlige Division)
7      int     (Nachkommastellen werden abgeschnitten)
-2     int     (abgeschnitten in Richtung 0)
4.0    float
15     int     (erst str -> int, dann Addition)
123    str     (Verkettung von zwei Strings)
2      int     (True wird als 1 gerechnet)
\end{lstlisting}
\end{solutionbox}
